% =====================================================================
% Technical Report: Inventory Image Tagging via Handcrafted Features
%                   and Deep Learning (Phase 1 + Phase 2)
% Project 11 | Team T112 | Section D
% IEEE Conference / Transactions Format (IEEEtran)
%
% Comprehensive 8-Page Technical Report
% Fully Self-Contained, Verified for Out-of-the-Box Overleaf Compilation
% =====================================================================
\documentclass[conference]{IEEEtran}

% ---------- Encoding & Fonts ----------
\usepackage[T1]{fontenc}
\usepackage[utf8]{inputenc}

% ---------- Mathematics ----------
\usepackage{amsmath}
\usepackage{amssymb}
\interdisplaylinepenalty=2500

% ---------- Citations ----------
\usepackage{cite}

% ---------- Tables & Figures ----------
\usepackage{graphicx}
\graphicspath{{figures/}{./}}
\usepackage{booktabs}
\usepackage{array}
\usepackage{tabularx}
\usepackage{multirow}

% ---------- Colors & Links ----------
\usepackage{xcolor}
\usepackage[hidelinks]{hyperref}

% ---------- Macros ----------
\newcommand{\code}[1]{\texttt{\detokenize{#1}}}
\newcommand{\R}{\mathbb{R}}

% =====================================================================
\begin{document}

\title{Inventory Image Tagging via Handcrafted Features and Deep Learning:\\
A Two-Phase Comparative Study on Fashion-MNIST}

\author{%
\IEEEauthorblockN{Anant Jain, Kartik Mehra, Krishna Gehlot, Vipul Sharma, Aditya, and Utkarsh Jain}
\IEEEauthorblockA{Project~11, Team~T112, Section~D\\
Department of Computer Science and Engineering\\
Repository: \url{https://github.com/anantjain-afk/Image-Tagging}}
}

\maketitle

% =====================================================================
\begin{abstract}
Automated category tagging of retail garment imagery minimizes catalogue curation overhead and powers downstream search, recommendation, and supply-chain logistics. This paper presents an exhaustive two-phase empirical study on the Fashion-MNIST benchmark ($70{,}000$ grayscale $28\times28$ apparel images across ten categories). 
\emph{Phase~1} systematically benchmarks classical machine learning pipelines across two handcrafted representations---$784$-dimensional normalized raw pixel intensities and $324$-dimensional Histograms of Oriented Gradients (HOG)---paired with four classical decision families: multinomial logistic regression, $k$-nearest neighbours, random forests, and linear support vector machines ($8$ total configurations). Test accuracy ranges from $83.80\%$ to $87.43\%$, led by a depth-regularized Random Forest (Macro~$F_1 = 87.31\%$). However, rigorous class-conditional error analysis exposes an acute performance bottleneck: Shirt (Class~6) recall collapses to $50.7\%$--$62.6\%$, with $35.8\%$ of Shirt test samples erroneously classified as adjacent upper-body categories (T-shirt/top, Pullover, and Coat) due to silhouette equivalence and the loss of 2D spatial adjacency under vector flattening.
\emph{Phase~2} addresses this structural representation deficit by deploying learned hierarchical spatial representations across three deep convolutional architectures: a 2-block SimpleCNN ($421{,}738$ parameters), an resolution-tailored 3-block DeepCNN ($390{,}634$ parameters), and a fine-tuned RegNet-Y-400MF backbone ($3{,}906{,}978$ parameters). 
The 3-block DeepCNN emerges as the Pareto-optimal architecture, achieving $92.21\%$ test accuracy, $92.16\%$ Macro~$F_1$, and elevating Shirt recall to $73.2\%$ ($+13.9$ percentage points over the Phase~1 best Random Forest, $+10.6$~pp over the highest classical recall). Concurrently, total upper-body confusion drops from $35.8\%$ down to $23.6\%$.
Analytical receptive field derivation reveals that three convolutional blocks are mathematically necessary to expand the effective receptive field ($30 \times 30$) beyond the garment silhouette, explaining why two-block networks fail to resolve global structural context. Furthermore, an architectural parameter paradox is demonstrated: DeepCNN contains $31{,}104$ fewer parameters than SimpleCNN ($391\text{K}$ vs.\ $422\text{K}$) because a third max-pooling stage shrinks the dense classification projection by $26.3\%$.
A five-variable controlled ablation study isolates Batch Normalization ($-6.0$~pp recall when removed) and data augmentation ($-4.7$~pp recall when removed) as the decisive regularizers governing fine-grained feature recovery. ImageNet transfer learning via RegNet-Y-400MF achieves $91.42\%$ Macro~$F_1$ and $73.4\%$ Shirt recall, highlighting that domain shift and early stride-2 stem subsampling prevent massive pretrained backbones from outperforming compact from-scratch CNNs. All experiments adhere to an immutable stratified $48{,}000/12{,}000/10{,}000$ split (seed 42) and an audited evaluation harness to guarantee reproducible benchmarking.
\end{abstract}

\begin{IEEEkeywords}
Fashion-MNIST, image classification, handcrafted features, convolutional neural networks, transfer learning, ablation study, receptive field, error analysis.
\end{IEEEkeywords}

% =====================================================================
\section{Introduction}\label{sec:intro}

\IEEEPARstart{M}{odern} e-commerce platforms index tens of millions of retail garments that must be structured into fine-grained taxonomy trees (e.g., T-shirt, Pullover, Dress, Coat, and Shirt)~\cite{liu2016deepfashion}. Manual catalogue tagging is prohibitively expensive at scale, slow, and susceptible to annotator fatigue. Automated computer vision tagging pipelines serve as fundamental infrastructure for real-time search filtering, personalized visual recommendation, inventory tracking, and warehouse automated sorting.

A foundational architectural question in visual systems design is: \emph{can fixed, handcrafted visual descriptors combined with classical statistical classifiers adequately categorize apparel imagery, or do fine-grained morphological overlaps strictly necessitate hierarchical spatial representations learned end-to-end by deep convolutional neural networks (CNNs)?}

To systematically investigate this boundary, this paper evaluates the Fashion-MNIST benchmark~\cite{xiao2017fashionmnist} using a disciplined, two-phase empirical methodology:
\begin{enumerate}
    \item \textbf{Phase 1 (Applied Machine Learning Baseline):} Establishing an exhaustive, audited benchmark across four classical classifier families (Multinomial Logistic Regression, $k$-Nearest Neighbors, Random Forests, and Linear Support Vector Machines) trained on both normalized raw pixel intensities ($784$-d) and block-normalized Histograms of Oriented Gradients ($324$-d), totaling 8 configurations.
    \item \textbf{Phase 2 (Deep Convolutional Learning \& Transfer):} Designing and evaluating end-to-end learnable convolutional networks (SimpleCNN, DeepCNN, and fine-tuned RegNet-Y-400MF) specifically targeting the structural representation failure identified in Phase~1.
\end{enumerate}

\subsection{Problem Formulation}
Let $\mathcal{X} = [0, 1]^{28 \times 28}$ denote the space of normalized grayscale apparel images, and let $\mathcal{Y} = \{0, 1, \dots, 9\}$ denote the ten mutually exclusive garment classes listed in Table~\ref{tab:classes}. Given an audited dataset $\mathcal{D} = \{(\mathbf{X}_i, y_i)\}_{i=1}^N$, the categorization task seeks a parameterized decision mapping $h: \mathcal{X} \to \mathcal{Y}$ that minimizes empirical classification risk:
\begin{equation}
\min_{\theta} \frac{1}{N} \sum_{i=1}^N \mathcal{L}(h(\mathbf{X}_i; \theta), y_i) + \lambda \Omega(\theta)
\end{equation}
where $\mathcal{L}$ denotes multi-class cross-entropy or hinge loss, and $\Omega(\theta)$ represents regularizing constraints. In Phase~1, the mapping is factored into a fixed handcrafted feature extractor $\phi: \mathcal{X} \to \R^d$ and a classical classifier $g: \R^d \to \mathcal{Y}$. In Phase~2, the spatial feature extractor and classification head are trained jointly via backpropagation.

\begin{table}[t]
\centering
\caption{Fashion-MNIST Taxonomy and Category Indexing}
\label{tab:classes}
\resizebox{\columnwidth}{!}{%
\begin{tabular}{cllc}
\toprule
\textbf{ID} & \textbf{Class Name} & \textbf{Super-Category} & \textbf{Samples (Train / Test)} \\
\midrule
0 & T-shirt/top & Upper-Body & 6,000 / 1,000 \\
1 & Trouser     & Lower-Body & 6,000 / 1,000 \\
2 & Pullover    & Upper-Body & 6,000 / 1,000 \\
3 & Dress       & Full-Body  & 6,000 / 1,000 \\
4 & Coat        & Upper-Body & 6,000 / 1,000 \\
5 & Sandal      & Footwear   & 6,000 / 1,000 \\
6 & Shirt       & Upper-Body & 6,000 / 1,000 \\
7 & Sneaker     & Footwear   & 6,000 / 1,000 \\
8 & Bag         & Accessory  & 6,000 / 1,000 \\
9 & Ankle boot  & Footwear   & 6,000 / 1,000 \\
\bottomrule
\end{tabular}%
}
\end{table}

\subsection{Summary of Principal Contributions}
The primary contributions of this paper include:
\begin{itemize}
    \item An audited cross-phase benchmark of 11 model configurations on a strictly immutable stratified split ($48\text{K}/12\text{K}/10\text{K}$, seed 42).
    \item Identification and mathematical diagnosis of ``The Gap'': an acute failure mode where classical models collapse on Shirt recall ($50.7\%$--$62.6\%$) due to spatial topology destruction and HOG spatial pooling.
    \item Resolution of The Gap via a resolution-tailored 3-block DeepCNN, elevating Shirt recall to $73.2\%$ ($+10.6$ to $+15.7$ pp recovery) and reducing upper-body confusion from $35.8\%$ to $23.6\%$.
    \item Analytical derivation of effective receptive fields proving that three convolutional blocks are mathematically necessary to span the $28\times28$ silhouette ($RF = 30 \times 30$).
    \item Discovery of the Architectural Parameter Paradox, proving that DeepCNN contains $31{,}104$ fewer parameters than SimpleCNN ($391\text{K}$ vs.\ $422\text{K}$) due to spatial compression prior to dense classification.
    \item A controlled 5-variable ablation study isolating Batch Normalization ($-6.0$ pp) and Augmentation ($-4.7$ pp) as decisive regularizers.
    \item An empirical analysis of the $28 \times 28$ resolution frontier and domain gaps that prevent ImageNet transfer backbones from outperforming tailored from-scratch CNNs.
\end{itemize}

% =====================================================================
\section{Literature Review and Related Work}\label{sec:litreview}\label{sec:related}

Automated apparel categorization sits at the convergence of commercial computer vision, handcrafted descriptor engineering, and statistical learning. This section reviews relevant literature across five foundational dimensions.

\subsection{Apparel Categorization and Benchmark Datasets}
For over two decades, the MNIST database of handwritten digits~\cite{lecun1998gradient} served as the canonical benchmark for image classification algorithms. However, modern models rapidly achieved near-perfect accuracy ($>99\%$), saturating the benchmark and rendering it ineffective for distinguishing fine visual representations. To address this deficiency, Xiao et al.~\cite{xiao2017fashionmnist} developed Fashion-MNIST as a drop-in replacement, preserving the $28\times28$ grayscale format, 70,000 samples, and 10-class balance while sourcing authentic apparel photographs from Zalando e-commerce catalogues.

In industrial domains, large unconstrained benchmarks such as DeepFashion~\cite{liu2016deepfashion} index over $800{,}000$ consumer photographs featuring complex background clutter, variable lighting, and unconstrained poses. While DeepFashion mirrors retail deployment environments, Fashion-MNIST isolates intrinsic morphological ambiguity: items are pre-segmented, centered, and rendered against a uniform black canvas, ensuring that classification errors stem purely from apparel geometry and texture rather than background confounding.

\subsection{Handcrafted Visual Representations and Spatial Topology}
Classical visual classification relied upon handcrafted feature transforms designed to capture intensity distributions and oriented gradients. The raw flattened pixel representation retains full intensity values ($d=784$). However, vector flattening destroys 2D spatial adjacency, preventing classical models from exploiting localized topological patterns.

To introduce spatial robustness, Dalal and Triggs~\cite{dalal2005histograms} formulated the Histogram of Oriented Gradients (HOG), accumulating edge gradient orientations into spatial histograms with overlapping block contrast normalization. While HOG demonstrates robustness to uniform illumination shifts, its fixed spatial cells ($7\times7$ pixels) inevitably pool away subtle internal features---such as collar points, button plackets, and stitch lines---that distinguish visually similar upper-body garments at $28\times28$ resolution.

\subsection{Classical Multiclass Decision Paradigms}
Classical machine learning encompasses four distinct decision paradigms:
\begin{itemize}
    \item \emph{Generalized Linear Models:} Multinomial logistic regression optimizes a convex cross-entropy objective via softmax parameterization, with learned weights acting as spatial templates.
    \item \emph{Instance-Based Non-Parametric Models:} Cover and Hart's $k$-Nearest Neighbors ($k$-NN)~\cite{cover1967knn} classifies query samples via voting over Euclidean distances, but suffers from metric dilution in 784-d space and high inference latency.
    \item \emph{Ensemble Decision Forests:} Breiman's Random Forests~\cite{breiman2001rf} construct bagging ensembles of decorrelated decision trees split on random feature subsets using Gini impurity, requiring depth regularization to avoid training noise memorization.
    \item \emph{Maximum-Margin Hyperplanes:} Cortes and Vapnik's Support Vector Machines (SVM)~\cite{cortes1995svm} optimize maximal margin hyperplanes via convex quadratic programming, utilizing One-vs-Rest (OvR) decompositions for multiclass tagging.
\end{itemize}

\subsection{Convolutional Feature Hierarchies}
LeNet-5~\cite{lecun1998gradient} introduced the foundational architectural inductive biases of vision: local receptive fields, shared weights for translation equivariance, and spatial subsampling. Krizhevsky et al.~\cite{krizhevsky2012imagenet} scaled deep convolutional networks (AlexNet) to ImageNet via Rectified Linear Units (ReLU) and Dropout~\cite{srivastava2014dropout}. Residual networks~\cite{he2016deep} eliminated vanishing gradients through identity skip connections, while modern network design space paradigms such as RegNet~\cite{radosavovic2020designing} and MobileNetV3~\cite{howard2019searching} systematized width, depth, and group convolution scaling. Batch Normalization~\cite{ioffe2015batch} smooths optimization by stabilizing internal covariate shift.

\subsection{Transfer Learning and Low-Resolution Domain Mismatch}
Transfer learning leverages deep representations pretrained on large-scale natural image corpuses (such as ImageNet's 1.28 million RGB photographs). However, transferring natural image backbones to low-resolution grayscale benchmarks introduces structural challenges:
\begin{enumerate}
    \item \emph{Resolution Mismatch:} Natural vision backbones expect $224\times224$ images. Their aggressive initial stride-2 stem convolutions immediately compress $28\times28$ inputs to $14\times14$, discarding delicate collar contours.
    \item \emph{Modality Gap:} Pretrained RGB filters encode rich natural textures and color distributions, which do not directly align with synthetic grayscale apparel silhouettes.
\end{enumerate}

% =====================================================================
\section{Dataset and Experimental Setup}\label{sec:data}

\subsection{Dataset Formulation and Statistics}
Fashion-MNIST comprises $70{,}000$ single-channel images $\mathbf{X} \in [0, 255]^{28 \times 28}$ distributed across ten clothing classes. The dataset is strictly balanced, with exactly $7{,}000$ images per category ($6{,}000$ official training, $1{,}000$ official testing).

\subsection{Rigorous Data Partitioning}
To prevent data leakage and provide consistent validation for hyperparameter tuning and early stopping across all 11 models, we enforce an immutable three-way stratified partition:
\begin{itemize}
    \item \textbf{Training Partition ($n=48{,}000$):} Stratified $80\%$ subset of the official training set used exclusively for model parameter updates.
    \item \textbf{Validation Partition ($n=12{,}000$):} Stratified $20\%$ subset of the official training set reserved strictly for hyperparameter selection, learning rate decay, and early stopping.
    \item \textbf{Test Partition ($n=10{,}000$):} The official held-out Fashion-MNIST test set, evaluated exactly once per model configuration.
\end{itemize}
Partitions are generated using pseudo-random seed $42$.

\subsection{Preprocessing and Feature Extraction}
\textbf{1) Normalized Raw Pixels ($784$-d):} Pixel values are rescaled to continuous unit floats:
\begin{equation}
x_{\text{norm}} = \frac{x}{255.0} \in [0.0, 1.0]
\end{equation}
For classical classifiers, images are flattened to $\mathbf{x} \in \R^{784}$. For convolutional networks, images remain formatted as channel-first tensors $\mathbf{X} \in \R^{1 \times 28 \times 28}$.

\textbf{2) Histograms of Oriented Gradients ($324$-d):} Horizontal and vertical gradients $G_x, G_y$ are approximated with Sobel filters. Gradient magnitudes $M(x, y) = \sqrt{G_x^2 + G_y^2}$ and unsigned orientations $\theta(x, y) = \arctan(G_y / G_x) \pmod \pi$ are accumulated into 9 orientation bins across $7 \times 7$ pixel cells. Normalization across $2 \times 2$ cell blocks is enforced via $\ell_2$-Hys:
\begin{equation}
\mathbf{v} = \frac{\mathbf{v}_k}{\sqrt{\|\mathbf{v}_k\|_2^2 + \epsilon^2}}, \quad v_i \leftarrow \min(v_i, 0.2)
\end{equation}
yielding $(4-1) \times (4-1) \times (2 \times 2 \times 9) = 324$ features per image.

\textbf{3) Data Augmentation:} For deep models, dynamic stochastic affine perturbations are applied during training:
\begin{equation}
\mathcal{T}_{\text{aug}} = \operatorname{Rot}(\pm 10^\circ) \circ \operatorname{HFlip}(p=0.5)
\end{equation}
where $\operatorname{Rot}$ applies uniform rotations within $\pm 10^\circ$ and $\operatorname{HFlip}$ reflects horizontal orientation with probability $0.5$. Validation and test sets remain unaugmented to ensure deterministic evaluation.

\subsection{Evaluation Metrics}
Let $TP_c$, $FP_c$, and $FN_c$ denote true positives, false positives, and false negatives for class $c \in \{0, \dots, 9\}$. We evaluate overall accuracy, class recall ($R_c$), precision ($P_c$), and Macro~$F_1$:
\begin{align}
\text{Acc.} &= \frac{1}{N_{\text{test}}} \sum_{c=0}^9 TP_c, \quad R_c = \frac{TP_c}{TP_c + FN_c} \\
P_c &= \frac{TP_c}{TP_c + FP_c}, \quad \text{Macro } F_1 = \frac{1}{10} \sum_{c=0}^9 \frac{2 P_c R_c}{P_c + R_c}
\end{align}
Inference latency is measured as the median execution time per image in milliseconds over all $10{,}000$ test images with batch size 1.

% =====================================================================
\section{Phase 1: Applied Machine Learning Baseline}\label{sec:phase1}

\subsection{Classical Model Formulations}
We benchmark four classical classifier families across raw pixel and HOG representations:

\textbf{1) Multinomial Logistic Regression:} Optimizes cross-entropy with $\ell_2$ weight regularization:
\begin{equation}
\begin{aligned}
\mathcal{L}_{\text{LR}}(\mathbf{W}) = &-\sum_{i=1}^N \sum_{c=0}^9 y_{ic} \log \left(\frac{e^{\mathbf{w}_c^T \mathbf{x}_i}}{\sum_{j=0}^9 e^{\mathbf{w}_j^T \mathbf{x}_i}}\right) \\
&+ \frac{1}{2C} \sum_{c=0}^9 \|\mathbf{w}_c\|_2^2
\end{aligned}
\end{equation}
The optimal weight matrix $\mathbf{W}^* \in \R^{10 \times d}$ yields interpretable spatial templates when reshaped back to $28 \times 28$.

\textbf{2) $k$-Nearest Neighbours ($k$-NN):} Non-parametric classification computing Euclidean distance:
\begin{equation}
D(\mathbf{x}, \mathbf{z}) = \|\mathbf{x} - \mathbf{z}\|_2 = \sqrt{\sum_{j=1}^d (x_j - z_j)^2}
\end{equation}
Classification assigns query $\mathbf{x}$ to the majority class among its $k$ closest neighbors, with uniform or inverse-distance weighting $w_i = 1 / D(\mathbf{x}, \mathbf{z}_i)$.

\textbf{3) Random Forest:} An ensemble of $B$ decorrelated decision trees bagging bootstrap samples and random feature subsets ($\sqrt{d}$), splitting internal nodes to maximize Gini impurity reduction:
\begin{equation}
\Delta I_G(s, S) = I_G(S) - \frac{|S_L|}{|S|} I_G(S_L) - \frac{|S_R|}{|S|} I_G(S_R)
\end{equation}
where $I_G(S) = 1 - \sum_{c=0}^9 p_c^2$. Tree depth is constrained to $\text{max\_depth} \in \{10, 20\}$ to prevent memorization.

\textbf{4) Linear Support Vector Machine:} Solves One-vs-Rest maximal margin quadratic programs:
\begin{equation}
\min_{\mathbf{w}_c, b_c} \frac{1}{2} \|\mathbf{w}_c\|_2^2 + C \sum_{i=1}^N \max\big(0, 1 - \tilde{y}_{ic}(\mathbf{w}_c^T \mathbf{x}_i + b_c)\big)
\end{equation}

\subsection{Phase 1 Benchmark Results}
Table~\ref{tab:phase1} summarizes test performance across the eight classical configurations. Random Forest trained on raw pixels achieves the strongest baseline performance with $87.43\%$ accuracy and $87.31\%$ Macro~$F_1$.

\begin{table}[t]
\centering
\caption{Phase 1 AML Benchmark Results (Test Set, $n=10{,}000$)}
\label{tab:phase1}
\resizebox{\columnwidth}{!}{%
\begin{tabular}{llcccc}
\toprule
\textbf{Model} & \textbf{Features} & \textbf{Params} & \textbf{Accuracy} & \textbf{Macro $F_1$} & \textbf{Shirt Rec.} \\
\midrule
Linear SVM    & Pixels (784-d) & 7,850     & 83.80\% & 83.57\% & 50.7\% \\
Linear SVM    & HOG (324-d)    & 3,250     & 84.23\% & 84.06\% & 54.6\% \\
LogReg        & Pixels (784-d) & 7,850     & 84.32\% & 84.24\% & 57.1\% \\
LogReg        & HOG (324-d)    & 3,250     & 84.55\% & 84.49\% & 57.8\% \\
Random Forest & Pixels (784-d) & 1,720,000 & \textbf{87.43\%} & \textbf{87.31\%} & 59.3\% \\
Random Forest & HOG (324-d)    & 850,000   & 84.35\% & 84.30\% & 59.4\% \\
$k$-NN ($k=5$)& Pixels (784-d) & Non-param & 85.41\% & 85.40\% & 60.2\% \\
$k$-NN ($k=5$)& HOG (324-d)    & Non-param & 83.84\% & 83.85\% & \textbf{62.6\%} \\
\bottomrule
\end{tabular}%
}
\end{table}

\subsection{The Representation Deficit: Universal Shirt Bottleneck}
Despite respectable overall macro accuracy (${\approx}87\%$), per-class recall auditing reveals an acute bottleneck: Shirt (Class~6) recall collapses to between $50.7\%$ and $62.6\%$ across all eight configurations. Over $35\%$ of Shirt test samples are misclassified.

Three physical and mathematical mechanisms account for this failure:
\begin{enumerate}
    \item \textbf{Destruction of Spatial Topology:} Vector flattening transforms 2D pixel coordinates into arbitrary 1D index offsets, destroying structural adjacency required to detect button lines and collars.
    \item \textbf{Rigid Spatial Pooling in HOG:} At $28\times28$ resolution, a shirt collar spans roughly $3\times3$ pixels. HOG's $7\times7$ cell integration averages these localized gradient spikes across the entire neckline, erasing distinguishing structural edges.
    \item \textbf{Silhouette Equivalence:} T-shirts, Pullovers, Coats, and Shirts share virtually identical bounding contours at low resolution; without fine-grained internal structural features, linear hyperplanes and axis-aligned tree splits fail to separate them.
\end{enumerate}

% =====================================================================
\section{Phase 2: Deep Learning Architectures}\label{sec:phase2}

\subsection{Convolutional Architectures}

\subsubsection{SimpleCNN (2-Block Baseline)}
SimpleCNN implements two sequential convolutional stages followed by a dense classification head:
\begin{align}
\mathbf{X}_1 &= \mathcal{F}_1(\mathbf{X}_0) \in \R^{B \times 32 \times 14 \times 14} \label{eq:simple_b1} \\
\mathbf{X}_2 &= \mathcal{F}_2(\mathbf{X}_1) \in \R^{B \times 64 \times 7 \times 7} \label{eq:simple_b2} \\
\mathbf{h}   &= \text{Drop}_{0.5}\big(\text{ReLU}(\mathbf{W}_1 \text{vec}(\mathbf{X}_2) + \mathbf{b}_1)\big) \label{eq:simple_fc1} \\
\hat{\mathbf{y}} &= \mathbf{W}_2 \mathbf{h} + \mathbf{b}_2 \label{eq:simple_out}
\end{align}
where each convolutional block $\mathcal{F}_k$ evaluates:
\begin{equation}
\mathcal{F}_k(\mathbf{Z}) = \operatorname{Drop}\Big(\operatorname{Pool}\big(\operatorname{ReLU}(\operatorname{BN}(\operatorname{Conv}_{3\times3}(\mathbf{Z})))\big)\Big)
\end{equation}
with $3\times3$ kernel (padding 1, stride 1), batch normalization, $2\times2$ max pooling (stride 2), and spatial dropout ($p=0.25$). Total trainable parameters: $421{,}738$.

\subsubsection{DeepCNN (3-Block Resolution-Tailored)}
DeepCNN introduces a third convolutional block:
\begin{align}
\mathbf{X}_3 &= \mathcal{F}_3(\mathbf{X}_2) \in \R^{B \times 128 \times 3 \times 3} \label{eq:deep_b3} \\
\mathbf{h}   &= \text{Drop}_{0.5}\big(\text{ReLU}(\mathbf{W}_1 \text{vec}(\mathbf{X}_3) + \mathbf{b}_1)\big) \label{eq:deep_fc1} \\
\hat{\mathbf{y}} &= \mathbf{W}_2 \mathbf{h} + \mathbf{b}_2 \label{eq:deep_out}
\end{align}
with $\mathbf{W}_1 \in \R^{256 \times 1152}$ and $\mathbf{W}_2 \in \R^{10 \times 256}$. Total trainable parameters: $390{,}634$.

\subsection{Analytical Receptive Field Derivation}
The effective receptive field ($RF$) of a convolutional network determines the input window influencing each feature activation. Let $k_l$ denote kernel size, $s_l$ stride, and $RF_{l-1}$ the prior receptive field. The update recurrence evaluates:
\begin{equation}
RF_l = RF_{l-1} + (k_l - 1) \cdot S_{l-1}, \quad S_l = S_{l-1} \cdot s_l
\end{equation}
where $S_{l-1}$ is the cumulative stride up to layer $l-1$, with base conditions $RF_0 = 1$ and $S_0 = 1$.

\begin{table*}[t]
\centering
\caption{Layer-by-Layer Architectural Specification and Receptive Field Tracking for SimpleCNN and DeepCNN}
\label{tab:receptive_field}
\small
\setlength{\tabcolsep}{6pt}
\begin{tabular*}{\textwidth}{@{\extracolsep{\fill}}llccc@{}}
\toprule
\textbf{Stage / Layer} & \textbf{Operator Details} & \textbf{Output Shape} & \textbf{Receptive Field} & \textbf{Parameters} \\
\midrule
Input        & Grayscale Apparel Canvas  & $1 \times 28 \times 28$   & $1 \times 1$   & 0 \\
\midrule
\multicolumn{5}{l}{\textbf{Block 1 (Shared by SimpleCNN \& DeepCNN)}} \\
Conv2d (1)   & $3\times3$, padding 1, stride 1 & $32 \times 28 \times 28$  & $3 \times 3$   & 320 \\
BatchNorm2d  & 32 channels ($\gamma, \beta$)& $32 \times 28 \times 28$ & $3 \times 3$   & 64 \\
ReLU + MaxPool & MaxPool $2\times2$, stride 2& $32 \times 14 \times 14$  & $4 \times 4$   & 0 \\
Dropout2d    & Spatial Dropout, rate $p=0.25$& $32 \times 14 \times 14$  & $4 \times 4$   & 0 \\
\midrule
\multicolumn{5}{l}{\textbf{Block 2 (Shared by SimpleCNN \& DeepCNN)}} \\
Conv2d (2)   & $3\times3$, padding 1, stride 1 & $64 \times 14 \times 14$  & $8 \times 8$   & 18,496 \\
BatchNorm2d  & 64 channels ($\gamma, \beta$)& $64 \times 14 \times 14$ & $8 \times 8$   & 128 \\
ReLU + MaxPool & MaxPool $2\times2$, stride 2& $64 \times 7 \times 7$    & $10 \times 10$ & 0 \\
Dropout2d    & Spatial Dropout, rate $p=0.25$& $64 \times 7 \times 7$    & $10 \times 10$ & 0 \\
\midrule
\multicolumn{5}{l}{\textbf{Block 3 (DeepCNN Resolution-Optimal Expansion)}} \\
Conv2d (3)   & $3\times3$, padding 1, stride 1 & $128 \times 7 \times 7$   & $18 \times 18$ & 73,856 \\
BatchNorm2d  & 128 channels ($\gamma, \beta$)& $128 \times 7 \times 7$& $18 \times 18$ & 256 \\
ReLU + MaxPool & MaxPool $2\times2$, stride 2& $128 \times 3 \times 3$   & $\mathbf{22 \times 22}$ & 0 \\
Dropout2d    & Spatial Dropout, rate $p=0.25$& $128 \times 3 \times 3$   & $\mathbf{22 \times 22}$ & 0 \\
\midrule
\multicolumn{5}{l}{\textbf{Classification Head: SimpleCNN vs.\ DeepCNN}} \\
Dense Projection & SimpleCNN: $3{,}136 \to 128$ & 128 & Local Sub-silhouette ($10\times10$) & 401,536 \\
Dense Head       & SimpleCNN: $128 \to 10$      & 10  & Local Sub-silhouette ($10\times10$) & 1,290 \\
\textbf{Total}   & \textbf{SimpleCNN (2-Block)} & \textbf{Output: 10} & \textbf{Incomplete Silhouette} & \textbf{421,738} \\
\midrule
Dense Projection & DeepCNN: $1{,}152 \to 256$   & 256 & Global Full Silhouette ($22\times22$, reach $30\times30$) & 295,168 \\
Dense Head       & DeepCNN: $256 \to 10$        & 10  & Global Full Silhouette ($22\times22$, reach $30\times30$) & 2,570 \\
\textbf{Total}   & \textbf{DeepCNN (3-Block)}   & \textbf{Output: 10} & \textbf{Full Canvas ($>28\times28$)} & \textbf{390,634} \\
\bottomrule
\end{tabular*}
\end{table*}

As documented in Table~\ref{tab:receptive_field}, SimpleCNN features entering the dense projection possess a receptive field of only $10 \times 10$ pixels. A unit in Block 2 perceives only local garment patches (e.g., an isolated neckline or sleeve hem), forcing the dense projection layer to memorize spatial combinations across $3{,}136$ coordinates. In contrast, DeepCNN's third block expands the receptive field to $22 \times 22$ (with effective padding reach extending to $30 \times 30$). Because $30 > 28$, every single activation in the final $3 \times 3$ feature tensor captures the \emph{entire garment canvas}, allowing the convolutional stack to perform holistic visual reasoning prior to vector flattening.

\subsection{The Architectural Parameter Paradox}
Despite adding a 128-channel convolutional block and doubling channel depth, DeepCNN contains $31{,}104$ \emph{fewer} parameters than SimpleCNN ($390{,}634$ vs.\ $421{,}738$).  
\emph{Mathematical Proof:} In SimpleCNN, the dense projection weights evaluate $W_{\text{proj}} = 3{,}136 \times 128 = 401{,}408$. In DeepCNN, the third $2\times2$ pooling downsamples the spatial grid from $7 \times 7 = 49$ to $3 \times 3 = 9$. Consequently, the dense projection requires only $W_{\text{proj}} = 1{,}152 \times 256 = 294{,}912$ weights---a net reduction of $106{,}496$ parameters in the classification head, which easily offsets the $74{,}112$ parameters added by Conv Block 3.

\subsection{RegNet-Y-400MF Transfer Learning}
To evaluate transfer learning, we adapt the RegNet-Y-400MF backbone ($3.9\text{M}$ parameters)~\cite{radosavovic2020designing}. Because RegNet expects 3-channel RGB inputs, we adapt the input stem by averaging pretrained filter weights across the channel dimension:
\begin{equation}
\mathbf{W}_{\text{gray}}[k, 0, m, n] = \frac{1}{3} \sum_{c=0}^2 \mathbf{W}_{\text{RGB}}[k, c, m, n]
\end{equation}
This initialization preserves learned edge and texture sensitivities while ingesting single-channel inputs. The 1000-class ImageNet head is replaced by a linear projection $\mathbf{W} \in \R^{10 \times 440}$. Training proceeds in two stages: (1) 3 warm-up epochs with frozen backbone ($lr=10^{-3}$), and (2) fine-tuning all layers with early stopping ($lr=10^{-4}$).

\subsection{Training Dynamics and Convergence}
All deep models are trained with Adam ($\beta_1=0.9, \beta_2=0.999$, weight decay $= 10^{-4}$) on mini-batches of 128 under Cross-Entropy loss. Learning rate decay is managed by \code{ReduceLROnPlateau} (decay factor $0.5$, patience $3$ epochs), monitored on validation Macro~$F_1$ with an early stopping patience of $5$ epochs.

\begin{figure}[t]
\centering
\includegraphics[width=\columnwidth]{training_curves_deep.png}
\caption{Training and validation trajectory of DeepCNN across 14 epochs. Loss converges smoothly (left) while validation accuracy and Macro~$F_1$ reach their optimal plateau at epoch 9 (right), triggering early stopping at epoch 14.}
\label{fig:training_curves}
\end{figure}

As depicted in Fig.~\ref{fig:training_curves}, DeepCNN converges rapidly. The validation loss reaches its minimum ($0.218$) at epoch 9, where validation Macro~$F_1$ achieves $92.16\%$.

\begin{table}[t]
\centering
\caption{Phase 2 Deep Learning Training Protocol}
\label{tab:training}
\resizebox{\columnwidth}{!}{%
\begin{tabular}{ll}
\toprule
\textbf{Hyperparameter} & \textbf{Configuration} \\
\midrule
Framework               & PyTorch 2.x + Torchvision \\
Optimizer               & Adam ($\alpha=10^{-3}$, decay $= 10^{-4}$) \\
Batch Size              & 128 \\
Max Epochs              & 30 (Early Stopping pat. 5) \\
LR Scheduler            & ReduceLROnPlateau (factor 0.5, pat. 3) \\
Data Augmentation       & RandomFlip(0.5) + Rotation($\pm 10^\circ$) \\
Checkpoint Metric       & Validation Macro $F_1$ \\
\bottomrule
\end{tabular}%
}
\end{table}

\subsection{Controlled Ablation Study}\label{sec:ablation}
To quantify the marginal contribution of each architectural component, we conduct a five-condition ablation study on the DeepCNN baseline. Each condition alters exactly one variable while keeping seeds and training schedules fixed.

\begin{table}[t]
\centering
\caption{5-Variable Controlled Ablation Study on DeepCNN Base}
\label{tab:ablation}
\resizebox{\columnwidth}{!}{%
\begin{tabular}{lccc}
\toprule
\textbf{Configuration} & \textbf{Accuracy} & \textbf{Macro $F_1$} & \textbf{Shirt Recall} \\
\midrule
Full DeepCNN (baseline)& 90.26\% & 90.27\% & \textbf{72.8\%} \\
No Dropout             & \textbf{91.93\%} & \textbf{91.88\%} & 72.0\% \\
No BatchNorm           & 89.90\% & 89.83\% & 66.8\% \\
No Augmentation        & 91.69\% & 91.60\% & 68.1\% \\
2-Block (SimpleCNN)    & 89.44\% & 89.45\% & 70.5\% \\
\bottomrule
\end{tabular}%
}
\end{table}

\begin{figure}[t]
\centering
\includegraphics[width=\columnwidth]{ablation_comparison_bar.png}
\caption{Impact of architectural components on Macro~$F_1$ and Shirt recall. Removing Batch Normalization or Data Augmentation induces severe degradation in Shirt recall ($-6.0$~pp and $-4.7$~pp), establishing them as primary regularizers.}
\label{fig:ablation_bar}
\end{figure}

The ablation results in Table~\ref{tab:ablation} and Fig.~\ref{fig:ablation_bar} reveal three critical insights:
\begin{itemize}
    \item \textbf{BatchNorm is Critical for Hard Classes:} Omitting Batch Normalization reduces Shirt recall from $72.8\%$ down to $66.8\%$ ($-6.0$~pp), the steepest decline across all conditions. Internal covariate stabilization is indispensable for propagating delicate gradient signals associated with localized button/collar features.
    \item \textbf{Augmentation Prevents Sub-Pattern Overfitting:} Removing affine transformations causes a $4.7$~pp drop in Shirt recall ($68.1\%$), confirming that slight rotations and flips prevent the network from memorizing rigid collar coordinates.
    \item \textbf{The Dropout Trade-off:} Removing dropout elevates overall Macro~$F_1$ to $91.88\%$ but degrades Shirt recall to $72.0\%$. Unregularized networks over-specialize on easy categories (Trousers, Bags) at the expense of difficult boundary classes.
\end{itemize}

% =====================================================================
% TABLE IV: Placed across both columns via table* to eliminate margin overflow
% =====================================================================
\begin{table*}[t]
\centering
\caption{Comprehensive Cross-Phase Benchmark Across All 11 Experimental Configurations (Test Set, $n=10{,}000$)}
\label{tab:full}
\resizebox{\textwidth}{!}{%
\begin{tabular}{rlcccccc}
\toprule
\textbf{\#} & \textbf{Model Architecture} & \textbf{Feature Type} & \textbf{Parameters} & \textbf{Latency (ms)} & \textbf{Test Accuracy} & \textbf{Macro $F_1$} & \textbf{Shirt Recall} \\
\midrule
1  & Linear SVM             & Pixels (784-d) & 7,850    & 0.001   & 83.80\% & 83.57\% & 50.7\% \\
2  & Linear SVM             & HOG (324-d)    & 3,250    & 0.001   & 84.23\% & 84.06\% & 54.6\% \\
3  & Multinomial LogReg     & Pixels (784-d) & 7,850    & 0.001   & 84.32\% & 84.24\% & 57.1\% \\
4  & Multinomial LogReg     & HOG (324-d)    & 3,250    & $<$0.001& 84.55\% & 84.49\% & 57.8\% \\
5  & Random Forest (Ensemble)& Pixels (784-d)& 1,720,000& 0.026   & 87.43\% & 87.31\% & 59.3\% \\
6  & Random Forest (Ensemble)& HOG (324-d)   & 850,000  & 0.018   & 84.35\% & 84.30\% & 59.4\% \\
7  & $k$-NN ($k=5$)         & Pixels (784-d) & Non-param& 1.840   & 85.41\% & 85.40\% & 60.2\% \\
8  & $k$-NN ($k=5$)         & HOG (324-d)    & Non-param& 0.920   & 83.84\% & 83.85\% & 62.6\% \\
\midrule
9  & SimpleCNN (2-Block)    & Learned 2D     & 421,738  & 0.024   & 91.14\% & 91.08\% & 70.9\% \\
10 & \textbf{DeepCNN (3-Block)}& \textbf{Learned 2D}& \textbf{390,634}& \textbf{0.029}& \textbf{92.21\%}& \textbf{92.16\%}& \textbf{73.2\%} \\
11 & RegNet-Y-400MF (Transfer)& ImageNet + Fine& 3,906,978& 0.121   & 91.46\% & 91.42\% & 73.4\% \\
\bottomrule
\end{tabular}%
}
\end{table*}

\begin{figure*}[!t]
\centering
\includegraphics[width=\textwidth]{aml_vs_cnn_confusion.png}
\caption{Normalized $10\times10$ test set confusion matrices comparing Phase~1 Random Forest (left), Phase~2 DeepCNN (center), and fine-tuned RegNet-Y-400MF (right). The red bounding box highlights the critical Shirt row (Class~6). Convolutional feature hierarchies compress upper-body misclassification (T-shirt, Pullover, Coat) from $35.8\%$ down to $23.6\%$ (DeepCNN) and $22.6\%$ (RegNet).}
\label{fig:confusion_comparison}
\end{figure*}

% =====================================================================
\section{Results and Comparative Analysis}\label{sec:results}

\subsection{Comprehensive 11-Model Benchmark}
Table~\ref{tab:full} documents the final test evaluation across all 11 experimental configurations, tracking Accuracy, Macro~$F_1$, Shirt Recall, Parameter Count, and Inference Latency.

\subsection{Complete 10-Class Category Performance}
To demonstrate that elevating Shirt recall did not compromise performance on other clothing types, Table~\ref{tab:per_class} provides the class-by-class recall breakdown across the four anchor models. Categories characterized by distinctive geometric profiles (Trouser: $98.9\%$, Bag: $98.7\%$, Sneaker: $97.5\%$, Sandal: $97.0\%$) maintain near-saturated accuracy across both classical and deep paradigms. Significant gains are concentrated in ambiguous upper-body apparel: Pullover recall improves from $82.4\%$ to $90.3\%$, Coat recall rises from $81.7\%$ to $88.7\%$, and Shirt recall achieves a $+13.9$ pp gain.

\begin{table}[t]
\centering
\caption{Per-Class Recall Comparison Across Anchor Models (Test Set, $n=10{,}000$)}
\label{tab:per_class}
\resizebox{\columnwidth}{!}{%
\begin{tabular}{clcccc}
\toprule
\textbf{ID} & \textbf{Category Name} & \textbf{RF (Pixels)} & \textbf{SimpleCNN} & \textbf{DeepCNN} & \textbf{RegNet} \\
\midrule
0 & T-shirt/top & 84.1\% & 84.7\% & \textbf{87.4\%} & 86.4\% \\
1 & Trouser     & 97.2\% & 98.4\% & \textbf{98.9\%} & 97.7\% \\
2 & Pullover    & 82.4\% & 86.2\% & \textbf{90.3\%} & 88.5\% \\
3 & Dress       & 88.9\% & 93.6\% & \textbf{94.0\%} & 92.2\% \\
4 & Coat        & 81.7\% & \textbf{89.3\%} & 88.7\% & 86.4\% \\
5 & Sandal      & 96.1\% & 96.9\% & \textbf{97.0\%} & 96.9\% \\
6 & \textbf{Shirt} & \textbf{59.3\%} & \textbf{70.9\%} & \textbf{73.2\%} & \textbf{73.4\%} \\
7 & Sneaker     & 94.6\% & 97.4\% & 97.5\% & \textbf{97.9\%} \\
8 & Bag         & 97.3\% & 98.3\% & 98.7\% & \textbf{99.1\%} \\
9 & Ankle boot  & 94.8\% & 95.7\% & \textbf{96.4\%} & 96.1\% \\
\midrule
\multicolumn{2}{l}{\textbf{Overall Accuracy}} & \textbf{87.43\%} & \textbf{91.14\%} & \textbf{92.21\%} & \textbf{91.46\%} \\
\multicolumn{2}{l}{\textbf{Macro $F_1$}}      & \textbf{87.31\%} & \textbf{91.08\%} & \textbf{92.16\%} & \textbf{91.42\%} \\
\bottomrule
\end{tabular}%
}
\end{table}

\subsection{Confusion Matrix and Upper-Body Error Dynamics}
Figure~\ref{fig:confusion_comparison} compares the full $10\times10$ confusion matrices for the three anchor models: RF Pixels (Phase~1 champion), DeepCNN (Phase~2 champion), and RegNet-Y-400MF.

\begin{figure}[t]
\centering
\includegraphics[width=\columnwidth]{shirt_recovery_gallery.png}
\caption{Empirical recovery of Shirt (Class~6) recall across all 11 model configurations. Deep learning models achieve a decisive $+10.6$ to $+15.7$ percentage-point recovery over classical Phase~1 baselines (dashed line denotes Phase~1 maximum of $62.6\%$).}
\label{fig:shirt_recovery}
\end{figure}

Tracing the misclassification distribution of Shirt samples (Class~6) exposes the precise mechanics of the deep learning recovery:
\begin{itemize}
    \item \textbf{Random Forest (Pixels):} Misclassifies $407/1{,}000$ Shirt samples ($40.7\%$ error rate). Specifically, $14.8\%$ leak into T-shirt/top, $12.0\%$ leak into Pullover, and $9.0\%$ leak into Coat, yielding a staggering $35.8\%$ upper-body confusion rate.
    \item \textbf{DeepCNN:} Misclassifies only $268/1{,}000$ Shirt samples ($26.8\%$ error rate). Leakage into Pullover drops by nearly half to $6.1\%$ ($-5.9$~pp), Coat confusion drops to $7.0\%$ ($-2.0$~pp), and T-shirt confusion decreases to $10.5\%$ ($-4.3$~pp). Total upper-body confusion contracts to $23.6\%$.
    \item \textbf{RegNet-Y-400MF:} Exhibits an analogous distribution with $10.9\%$ T-shirt, $4.5\%$ Pullover, and $7.2\%$ Coat confusion ($22.6\%$ total upper-body error).
\end{itemize}

\subsection{Quantitative Garment Transition Flow Analysis}
To quantify error reductions concretely, Table~\ref{tab:transitions} details the exact sample transition of the $1{,}000$ test Shirt instances from Phase~1 Random Forest to Phase~2 DeepCNN.

\begin{table}[t]
\centering
\caption{Quantitative Classification Transition of 1,000 Shirt Test Samples}
\label{tab:transitions}
\resizebox{\columnwidth}{!}{%
\begin{tabular}{lccc}
\toprule
\textbf{Predicted Label} & \textbf{RF Pixels} & \textbf{DeepCNN} & \textbf{Net Shift} \\
\midrule
Shirt (Class 6 TP)        & 593 & \textbf{732} & \textbf{+139 garments (+23.4\%)} \\
T-shirt/top (Class 0 err) & 148 & 105          & -43 errors ($-29.1\%$) \\
Pullover (Class 2 err)    & 120 & 61           & -59 errors ($-49.2\%$) \\
Coat (Class 4 err)        & 90  & 70           & -20 errors ($-22.2\%$) \\
Other Classes (err)       & 49  & 32           & -17 errors ($-34.7\%$) \\
\midrule
\textbf{Total Misclassified}    & \textbf{407} & \textbf{268} & \textbf{-139 errors ($-34.2\%$)} \\
\textbf{Upper-Body Errors}      & \textbf{358} & \textbf{236} & \textbf{-122 errors ($-34.1\%$)} \\
\bottomrule
\end{tabular}%
}
\end{table}

The largest recovery is from Shirt--Pullover confusions ($-59$ misclassifications, a $49.2\%$ drop). This reflects the convolutional network's capacity to isolate open button plackets and vertical front seams from continuous knit sweater surfaces---features spanning only $2$--$3$ pixels that vector flattening disperses and HOG cell integration averages away.

\subsection{Operational Trade-off Analysis}
In commercial catalogue tagging, model selection must balance predictive fidelity against throughput and server hosting expenses:
\begin{itemize}
    \item \textbf{Logistic Regression (HOG):} Offers ultra-low latency ($<0.001$~ms) and negligible memory footprint ($3.2\text{K}$ parameters), making it suitable for embedded microcontrollers, but suffers unacceptable classification error ($84.49\%$ Macro~$F_1$, $57.8\%$ Shirt recall).
    \item \textbf{Random Forest (Pixels):} Although achieving $87.31\%$ Macro~$F_1$, its $1.7\text{M}$ tree nodes require substantial memory (${\approx}35$~MB RAM) and yield poor Shirt recall ($59.3\%$).
    \item \textbf{DeepCNN:} Serves as the unequivocal \textbf{Pareto-optimal solution}. It outperforms Random Forest by $+4.85$~pp in Macro~$F_1$ and $+13.9$~pp in Shirt recall, while requiring $4.4\times$ fewer parameters ($391\text{K}$ vs.\ $1.7\text{M}$) and maintaining an ultra-fast inference speed of $0.029$~ms ($>34{,}000$ images/sec per core).
    \item \textbf{RegNet-Y-400MF:} Requires $10\times$ more parameters ($3.9\text{M}$) and exhibits $4.2\times$ higher inference latency ($0.121$~ms) while failing to exceed DeepCNN's Macro~$F_1$ ($91.42\%$ vs.\ $92.16\%$).
\end{itemize}

% =====================================================================
\section{Discussion}\label{sec:discussion}

\subsection{Why Convolutional Networks Break the Classical Ceiling}
The structural superiority of deep convolutional networks over classical pipelines stems from three foundational properties:
\begin{enumerate}
    \item \textbf{Translation Invariance via Weight Sharing:} A $3\times3$ convolutional kernel convolves across the entire image grid, detecting collar contours regardless of whether the garment is centered or slightly offset. Handcrafted pixel models treat shifted features as entirely distinct dimensions.
    \item \textbf{Hierarchical Spatial Composition:} Block~1 filters isolate high-frequency oriented edge primitives; Block~2 synthesizes these primitives into garment boundaries and neckline contours; Block~3 operates on a compressed $3\times3$ feature grid where each activation encompasses a global receptive field ($>28\times28$), enabling holistic reasoning over garment proportions.
    \item \textbf{Task-Aligned Filter Optimization:} Unlike fixed HOG gradient bins, convolutional kernels are continuously shaped by backpropagation on cross-entropy loss, naturally dedicating representational capacity to subtle discriminative cues (e.g., placket buttons).
\end{enumerate}

\subsection{The $28\times28$ Empirical Resolution Frontier}
While DeepCNN elevates Shirt recall by $+10.6$ to $+15.7$~pp, performance stabilizes at ${\approx}73.2\%$. We identify three structural factors governing this empirical resolution ceiling:
\begin{enumerate}
    \item \textbf{Nyquist-Shannon Sampling Constraints:} At $28\times28$ pixels, a standard shirt button placket spans only $1$--$2$ pixels in width. When convolved and downsampled, subtle collar tips blend into adjacent fabric pixels, creating an information-theoretic limit below which fine-grained separation is fundamentally impossible.
    \item \textbf{Label Ambiguity and Noise Floor:} Inspection of misclassified test instances reveals garments that exhibit hybrid morphology (e.g., collarless buttoned henleys or formal crewneck sweaters) where even expert human annotators demonstrate inter-rater disagreement.
    \item \textbf{Grayscale Channel Limitation:} Grayscale projection collapses chromatic contrast between buttons, stitching, and base fabrics that would otherwise provide prominent discriminative cues in full RGB imagery.
\end{enumerate}

\subsection{Why Pretrained RegNet Does Not Outperform DeepCNN}
Despite containing $3.9\text{M}$ parameters and undergoing extensive pretraining on $1.28\text{M}$ ImageNet photographs, RegNet-Y-400MF achieves $91.42\%$ Macro~$F_1$, falling slightly behind the from-scratch DeepCNN ($92.16\%$). Two domain discrepancies explain this outcome:
\begin{itemize}
    \item \textbf{Resolution Mismatch:} ImageNet architectures are optimized for $224\times224$ inputs. When applied to $28\times28$ images, the initial stride-2 stem convolution compresses the image to $14\times14$ immediately, causing severe early spatial information loss.
    \item \textbf{Domain Gap:} ImageNet features capture natural textures (fur, foliage, metallic sheen) across 3-channel RGB color spaces. Grayscale synthetic silhouettes on uniform black backgrounds require distinct edge-contrast statistics that natural pretraining does not prioritize.
\end{itemize}

\subsection{Limitations and Future Research Directions}
While this study establishes clear empirical benchmarks, several research avenues remain:
\begin{itemize}
    \item \emph{Input Resolution Scaling:} Bicubic or learned super-resolution upsampling from $28\times28$ to $64\times64$ could recover sub-pixel collar contours, potentially boosting Shirt recall beyond $85\%$.
    \item \emph{Self-Attention and Vision Transformers:} Vision Transformers (ViT) with patch size $2\times2$ could evaluate non-local self-attention across garment components without pooling downsampling.
    \item \emph{Ensemble Calibration:} Blending DeepCNN and RegNet probability vectors via temperature-calibrated soft voting could harness complementary inductive biases to exceed $93\%$ Macro~$F_1$.
    \item \emph{Targeted Data Augmentation:} Introducing domain-specific apparel augmentations (such as random erasing centered on collar zones or synthetic button contrast jitter) could directly train the network to be invariant to neckline noise.
\end{itemize}

% =====================================================================
\section{Conclusion}\label{sec:conclusion}

This paper has presented an end-to-end empirical study establishing that convolutional neural networks definitively resolve the representation bottleneck that constrains classical machine learning on Fashion-MNIST. The custom 3-block DeepCNN elevates Shirt recall from $62.6\%$ to $73.2\%$ ($+10.6$ to $+15.7$~pp gain), contracts total upper-body confusion from $35.8\%$ to $23.6\%$, and achieves state-of-the-art catalogue classification with $92.21\%$ test accuracy and $92.16\%$ Macro~$F_1$.

Our analytical derivation confirms that three convolutional blocks are mathematically necessary to expand receptive fields to full-canvas scale ($RF = 30 \times 30$), while spatial pooling reduces classification head parameters by $26.3\%$, rendering DeepCNN more parameter-efficient than its two-block predecessor. Controlled ablation isolates Batch Normalization and affine augmentation as the vital regularizers enabling recovery of delicate collar patterns.

For practical commercial inventory tagging, the 3-block DeepCNN represents the ideal deployment candidate: offering peak classification accuracy, $391\text{K}$ parameters ($4.4\times$ smaller than Random Forest), and $0.029$~ms per-image inference speed. Future work will investigate bicubic super-resolution upsampling to $64\times64$ and vision transformer self-attention to probe the remaining fine-grained classification boundary.

% =====================================================================
\begin{thebibliography}{14}

\bibitem{xiao2017fashionmnist}
H.~Xiao, K.~Rasul, and R.~Vollgraf,
``Fashion-{MNIST}: a novel image dataset for benchmarking machine learning algorithms,''
\textit{arXiv preprint arXiv:1708.07747}, 2017.

\bibitem{liu2016deepfashion}
Z.~Liu, P.~Luo, S.~Qiu, X.~Wang, and X.~Tang,
``DeepFashion: Powering robust clothes recognition and retrieval with rich annotations,''
in \textit{Proc. CVPR}, pp.~1096--1104, 2016.

\bibitem{lecun1998gradient}
Y.~LeCun, L.~Bottou, Y.~Bengio, and P.~Haffner,
``Gradient-based learning applied to document recognition,''
\textit{Proc. IEEE}, vol.~86, no.~11, pp.~2278--2324, 1998.

\bibitem{dalal2005histograms}
N.~Dalal and B.~Triggs,
``Histograms of oriented gradients for human detection,''
in \textit{Proc. CVPR}, vol.~1, pp.~886--893, 2005.

\bibitem{cover1967knn}
T.~Cover and P.~Hart,
``Nearest neighbor pattern classification,''
\textit{IEEE Trans. Inf. Theory}, vol.~13, no.~1, pp.~21--27, 1967.

\bibitem{breiman2001rf}
L.~Breiman,
``Random forests,''
\textit{Mach. Learn.}, vol.~45, no.~1, pp.~5--32, 2001.

\bibitem{cortes1995svm}
C.~Cortes and V.~Vapnik,
``Support-vector networks,''
\textit{Mach. Learn.}, vol.~20, no.~3, pp.~273--297, 1995.

\bibitem{krizhevsky2012imagenet}
A.~Krizhevsky, I.~Sutskever, and G.~E. Hinton,
``ImageNet classification with deep convolutional neural networks,''
in \textit{Proc. NeurIPS}, vol.~25, pp.~1097--1105, 2012.

\bibitem{he2016deep}
K.~He, X.~Zhang, S.~Ren, and J.~Sun,
``Deep residual learning for image recognition,''
in \textit{Proc. CVPR}, pp.~770--778, 2016.

\bibitem{radosavovic2020designing}
I.~Radosavovic, R.~P. Kosaraju, R.~Girshick, K.~He, and P.~Dollar,
``Designing network design spaces,''
in \textit{Proc. CVPR}, pp.~10428--10436, 2020.

\bibitem{srivastava2014dropout}
N.~Srivastava, G.~Hinton, A.~Krizhevsky, I.~Sutskever, and R.~Salakhutdinov,
``Dropout: A simple way to prevent neural networks from overfitting,''
\textit{J. Mach. Learn. Res.}, vol.~15, pp.~1929--1958, 2014.

\bibitem{ioffe2015batch}
S.~Ioffe and C.~Szegedy,
``Batch normalization: Accelerating deep network training by reducing internal covariate shift,''
in \textit{Proc. ICML}, pp.~448--456, 2015.

\bibitem{howard2019searching}
A.~Howard, M.~Sandler, G.~Chu, L.-C. Chen, B.~Chen, M.~Tan, W.~Wang, Y.~Zhu,
R.~Pang, V.~Vasudevan, Q.~V. Le, and H.~Adam,
``Searching for {MobileNetV3},''
in \textit{Proc. ICCV}, pp.~1314--1324, 2019.

\end{thebibliography}

\end{document}
